% Generado por scripts/validar_modelo_financiero.py — no editar a mano.
\begin{table}[t]
  \centering
  \footnotesize\setlength{\tabcolsep}{3.5pt}
  \caption{Flujo de caja del escenario optimista, en USD. El flujo neto resta de los ingresos los costos, el impuesto a las ganancias y la inversión en capital de trabajo; en el año 5 esa inversión es negativa porque se recupera. Fuente: elaboración propia sobre la planilla del modelo financiero.}
  \label{tab:financiero-optimista-flujo}
  \begin{tabular}{c r r r r r r r}
    \hline
    \textbf{Año} & \textbf{Ingresos} & \textbf{Costos} & \textbf{Impuesto} & \textbf{Cap. de trabajo} & \textbf{Flujo neto} & \textbf{Descontado} & \textbf{Acumulado} \\
    \hline
    0 & 0,00 & 1.729,00 & 0,00 & 0,00 & (1.729,00) & (1.729,00) & (1.729,00) \\
    1 & 255.000,00 & 47.653,24 & 71.966,22 & 42.500,00 & 92.880,54 & 74.304,43 & 91.151,54 \\
    2 & 473.400,00 & 74.669,45 & 139.555,69 & 36.400,00 & 222.774,85 & 142.575,91 & 313.926,40 \\
    3 & 791.400,00 & 104.787,09 & 240.314,52 & 53.000,00 & 393.298,39 & 201.368,78 & 707.224,79 \\
    4 & 849.000,00 & 140.363,93 & 248.022,63 & 9.600,00 & 451.013,45 & 184.735,11 & 1.158.238,23 \\
    5 & 1.045.800,00 & 173.915,21 & 305.159,68 & (141.500,00) & 708.225,11 & 232.071,21 & 1.866.463,35 \\
    \hline
  \end{tabular}
\end{table}
